\ChapterImagePrelim[cap:resumen]{Resumen}{./images/fondo.png}\label{cap:resumen}
\mbox{}\\
\noindent
El presente trabajo de grado aborda el desarrollo de una notación unificada de modelado para la Infraestructura de Tecnologías de la Información (\NDITI), orientada a resolver las limitaciones de interoperabilidad y uniformidad conceptual presentes en las notaciones existentes. Este trabajo responde a la necesidad institucional de la Universidad del Quindío de fortalecer los procesos de enseñanza, aprendizaje e investigación en el campo de la infraestructura de \TI~, mediante la consolidación de un marco conceptual estandarizado que facilite la representación coherente de arquitecturas tecnológicas.
\\\\
Metodológicamente, el proyecto comprende: (\textbf{a}) reconocimiento de necesidades en torno al modelado de infraestructura de \TI~, (\textbf{b}) un estudio de mapeo sistemático (\SMS) para la identificación de la literatura relacionada a los lenguajes de modelado de Infraestructura de TI con el fin de propender por la toma de decisiones informadas, (\textbf{c}) la aplicación de la metodología de Análisis de Decisiones y Resolución (\DAR) del modelo \CMMI. Los resultados señalan a la notaciones \textbf{N1} como las más adecuada a utilizar. Finalmente, se desarrolla una propuesta de notación para la especificación de diseños de Infraestructura de \TI~utilizando como base el lenguaje de modelado <L1>.  El desarrollo se complementa con la elaboración de ejemplos y prototipos de modelado que ilustran la aplicabilidad de la \NDITI~en escenarios académicos y prácticos, permitiendo validar su claridad, expresividad y alineación con las necesidades formativas del programa. Considerando que la propuesta busca trascender las necesidades específicas del contexto institucional, se proyecta que la \NDITI~contribuya a la estandarización y mejora de los procesos de enseñanza, aprendizaje e investigación relacionados con el diseño y gestión de infraestructuras tecnológicas en entornos universitarios.





\ChapterImagePrelim[cap:abstract]{Abstract}{./images/fondo.png}\label{cap:abstract}
\mbox{}\\
\noindent
This thesis addresses the development of a unified modeling notation for Information Technology Infrastructure (\NDITI), aimed at overcoming interoperability and conceptual uniformity limitations found in existing notations. The work responds to the institutional need of the University of Quindío to strengthen teaching, learning, and research processes in the field of \IT infrastructure by establishing a standardized conceptual framework that facilitates the coherent representation of technological architectures.
\\\\
Methodologically, the project includes: (\textbf{a}) identifying needs related to \IT~infrastructure modeling, (\textbf{b}) conducting a systematic mapping study (\SMS) to identify related literature on IT infrastructure modeling languages in order to support informed decision-making, and (\textbf{c}) applying the Decision Analysis and Resolution (\DAR) methodology from the CMMI model. The results indicate that the \textbf{N1} notation is the most suitable for use. Finally, a notation proposal is developed for specifying \IT~infrastructure designs, based on the modeling language L1. The development is complemented by examples and modeling prototypes that demonstrate the applicability of \NDITI~in academic and practical scenarios, allowing the validation of its clarity, expressiveness, and alignment with the educational needs of the program. Considering that the proposal aims to go beyond the specific needs of the institutional context, \NDITI~is expected to contribute to the standardization and improvement of teaching, learning, and research processes related to the design and management of technological infrastructures in university environments.